\documentclass[11pt]{article}
\usepackage[utf8]{inputenc}
\usepackage[margin=1in]{geometry}
\usepackage{amsmath}
\usepackage{graphicx}
\usepackage{booktabs}
\usepackage{hyperref}
\usepackage{natbib}
\usepackage{lineno}
\usepackage{xcolor}
\usepackage{multirow}
\usepackage{float}

\linenumbers

\title{Antipsychotic Drugs Selectively Decorrelate Long-Range Interactions in Deep Cortical Layers}

\author{
Author~One$^{1,\dagger}$, Author~Two$^{1,\dagger}$, Author~Three$^{1}$, Author~Four$^{1}$, \\
Author~Five$^{1,*}$ \\[6pt]
$^{1}$Friedrich Miescher Institute for Biomedical Research, Basel, Switzerland \\[4pt]
$^{\dagger}$These authors contributed equally. \\
$^{*}$Corresponding author: corresponding@example.edu
}

\date{}

\begin{document}
\maketitle

\begin{abstract}
Antipsychotic drugs are the frontline treatment for psychosis, yet their cortical circuit-level effects remain poorly understood. Using widefield calcium imaging through crystal skull cranial windows in head-fixed mice across eight Cre lines, we show that layer 5 intratelencephalic (L5 IT, Tlx3-Cre) neurons most strongly distinguish closed-loop from open-loop visuomotor conditions. Clozapine selectively disrupted visuomotor integration in L5 IT neurons while leaving superficial layer responses largely intact. All three antipsychotic drugs tested---clozapine, aripiprazole, and haloperidol---strongly reduced L5 IT inter-areal correlations (mean $\Delta r = -0.31 \pm 0.04$, $p < 0.001$, hierarchical bootstrap), whereas the psychostimulant amphetamine did not ($\Delta r = -0.02 \pm 0.03$, $p = 0.58$). Antipsychotic treatment preferentially reduced negative prediction error (mismatch) responses and their propagation from V1 to V2am in L5 IT neurons ($\Delta F/F$ reduction: $42 \pm 8\%$, $p < 0.01$). These findings reveal that antipsychotics converge on a layer-specific mechanism---decorrelation of deep-layer long-range cortical communication---that may underlie their therapeutic effects in psychosis.

\vspace{0.5em}
\noindent\textbf{Keywords:} antipsychotic, cortical layers, visuomotor integration, prediction error, inter-areal correlations, clozapine
\end{abstract}

\section{Introduction}

Psychosis is characterized by an impaired ability to distinguish between internally generated and externally driven neural activity \citep{adams2013computational, friston2005theory}. Antipsychotic drugs, which remain the primary pharmacological treatment for psychotic disorders \citep{lieberman2018effectiveness, correll2017current}, have broad neuromodulator receptor binding profiles spanning dopamine, serotonin, histamine, and muscarinic receptors \citep{kapur2005dopamine, meltzer2013serotonin, seeman2006targeting}. Despite decades of clinical use, the cortical circuit-level effects of these drugs remain unclear.

The predictive processing framework provides a computational account of cortical function in which superficial layers (L2/3) signal prediction errors while deep layers (L5/L6) maintain internal models and generate predictions \citep{rao1999predictive, keller2018predictive}. In mouse primary visual cortex (V1), visuomotor mismatch responses---a form of negative prediction error arising when locomotion-driven predictions do not match visual feedback---have been localized to specific cortical layers and propagate to higher visual areas \citep{zmarz2016mismatch, jordan2020opposing, leinweber2017sensorimotor}. If antipsychotic drugs restore the distinction between self-generated and external signals, they might selectively affect prediction error signaling in specific cortical layers.

Here we used widefield calcium imaging through crystal skull cranial windows \citep{kim2016crystal} in head-fixed mice running on a spherical treadmill in a virtual corridor to systematically characterize the effects of antipsychotic drugs across cortical cell types and layers. Using eight transgenic Cre lines targeting distinct excitatory and inhibitory neuron populations \citep{daigle2018suite}, we mapped drug effects with layer and cell-type specificity. We tested three antipsychotic drugs with distinct receptor profiles---clozapine (atypical), aripiprazole (atypical, D2 partial agonist), and haloperidol (typical)---and compared them to the psychostimulant amphetamine.

\section{Results}

\subsection{Layer 5 IT neurons most strongly distinguish closed- from open-loop conditions}

We first established that different cortical layers carry distinct visuomotor signals by comparing calcium responses during closed-loop (locomotion drives visual flow) and open-loop (visual replay) locomotion onsets across Cre lines (Table~\ref{tab:cl_ol}).

\begin{table}[H]
\centering
\caption{Discrimination between closed-loop and open-loop locomotion onsets by cell type. Area under the ROC curve (AUC) for classifying closed vs.\ open-loop from calcium responses in V1. Hierarchical bootstrap 90\% CI. $n$ = number of mice.}
\label{tab:cl_ol}
\begin{tabular}{llcccc}
\toprule
\textbf{Cre Line} & \textbf{Layer/Type} & \textbf{$n$} & \textbf{AUC} & \textbf{90\% CI} & \textbf{Rank} \\
\midrule
Tlx3 & L5 IT & 25 & $\mathbf{0.84}$ & $[0.79, 0.89]$ & 1 \\
Ntsr1 & L6 & 3 & $0.79$ & $[0.68, 0.88]$ & 2 \\
SST & Interneuron & 5 & $0.71$ & $[0.63, 0.78]$ & 3 \\
VIP & Interneuron & 6 & $0.68$ & $[0.59, 0.76]$ & 4 \\
PV & Interneuron & 2 & $0.65$ & $[0.54, 0.75]$ & 5 \\
Scnn1a & L4 & 7 & $0.58$ & $[0.51, 0.65]$ & 6 \\
Cux2 & L2/3--L4 & 4 & $0.56$ & $[0.48, 0.63]$ & 7 \\
Emx1 & Pan-excitatory & 4 & $0.62$ & $[0.55, 0.69]$ & --- \\
\bottomrule
\end{tabular}
\end{table}

Tlx3-positive L5 IT neurons showed the strongest discrimination (AUC = 0.84), substantially exceeding superficial layer neurons (Cux2 AUC = 0.56; Scnn1a AUC = 0.58). This establishes L5 IT as the cortical population most sensitive to visuomotor coupling.

\subsection{Clozapine effects on inter-regional correlations are layer-specific}

We next tested whether clozapine (atypical antipsychotic) affected cortical correlations differently across layers. Inter-regional correlations were computed across 12 dorsal cortical ROIs before and after single-dose clozapine administration (Table~\ref{tab:cloz_corr}).

\begin{table}[H]
\centering
\caption{Change in mean inter-regional correlation ($\Delta r$) after clozapine administration by cell type. Hierarchical bootstrap with 90\% CI \citep{efron1993bootstrap}. Significant reductions at 90\% level are bolded.}
\label{tab:cloz_corr}
\begin{tabular}{llcccc}
\toprule
\textbf{Cre Line} & \textbf{Layer} & \textbf{$n$} & \textbf{$\Delta r$ (mean)} & \textbf{90\% CI} & \textbf{Effect Size ($d$)} \\
\midrule
Tlx3 & L5 IT & 5 & $\mathbf{-0.34}$ & $[-0.42, -0.26]$ & $2.14$ \\
C57BL/6 & Brain-wide & 4 & $\mathbf{-0.18}$ & $[-0.25, -0.11]$ & $1.28$ \\
Scnn1a & L4 & 7 & $-0.04$ & $[-0.11, 0.03]$ & $0.31$ \\
Cux2 & L2/3--L4 & 4 & $-0.03$ & $[-0.09, 0.04]$ & $0.22$ \\
\bottomrule
\end{tabular}
\end{table}

Clozapine dramatically reduced inter-regional correlations in L5 IT neurons ($\Delta r = -0.34$, 90\% CI $[-0.42, -0.26]$) while producing negligible effects in superficial layers (Scnn1a $\Delta r = -0.04$; Cux2 $\Delta r = -0.03$).

\subsection{All antipsychotics decorrelate L5 IT; amphetamine does not}

We compared the effects of three mechanistically distinct antipsychotics and the psychostimulant amphetamine on L5 IT inter-areal correlations (Table~\ref{tab:drug_comparison}).

\begin{table}[H]
\centering
\caption{Drug effects on L5 IT (Tlx3) distance-dependent inter-areal correlations. $\Delta r$ = change in mean correlation. Hierarchical bootstrap 90\% CI. $^{***}p < 0.001$; NS = not significant.}
\label{tab:drug_comparison}
\begin{tabular}{llccccc}
\toprule
\textbf{Drug} & \textbf{Class} & \textbf{$n$} & \textbf{$\Delta r$} & \textbf{90\% CI} & \textbf{$d$} & \textbf{Sig.} \\
\midrule
Clozapine & Atypical AP & 5 & $\mathbf{-0.34}$ & $[-0.42, -0.26]$ & $2.14$ & $^{***}$ \\
Aripiprazole & Atypical AP (D2pa) & 4 & $\mathbf{-0.31}$ & $[-0.40, -0.22]$ & $1.94$ & $^{***}$ \\
Haloperidol & Typical AP & 3 & $\mathbf{-0.28}$ & $[-0.38, -0.18]$ & $1.75$ & $^{***}$ \\
Amphetamine & Psychostimulant & 4 & $-0.02$ & $[-0.08, 0.04]$ & $0.13$ & NS \\
\bottomrule
\end{tabular}
\end{table}

All three antipsychotics produced strong L5 IT decorrelation ($\Delta r$ range: $-0.28$ to $-0.34$) with large effect sizes ($d > 1.7$). Amphetamine, despite increasing dopamine release, produced no significant decorrelation ($\Delta r = -0.02$, $p = 0.58$).

\subsection{Visuomotor prediction error responses}

We analyzed mismatch (negative prediction error) and drifting grating responses in V1 before and after drug treatment (Table~\ref{tab:pe}).

\begin{table}[H]
\centering
\caption{Effect of antipsychotic treatment on V1 mismatch and grating responses in Tlx3 (L5 IT) neurons. $\Delta F/F$ peak amplitude as percentage of pre-drug baseline. Hierarchical bootstrap 90\% CI.}
\label{tab:pe}
\begin{tabular}{llcccc}
\toprule
\textbf{Stimulus} & \textbf{Region} & \textbf{Pre-Drug (\%)} & \textbf{Post-AP (\%)} & \textbf{Reduction (\%)} & \textbf{90\% CI} \\
\midrule
Mismatch & V1 & $100$ & $58 \pm 8$ & $\mathbf{42 \pm 8}$ & $[34, 50]$ \\
Mismatch & V2am & $100$ & $51 \pm 11$ & $\mathbf{49 \pm 11}$ & $[38, 60]$ \\
Grating & V1 & $100$ & $84 \pm 9$ & $16 \pm 9$ & $[7, 25]$ \\
Grating & V2am & $100$ & $88 \pm 10$ & $12 \pm 10$ & $[2, 22]$ \\
\bottomrule
\end{tabular}
\end{table}

Antipsychotic treatment preferentially reduced mismatch responses (42\% reduction in V1, 49\% in V2am) compared to grating responses (16\% in V1, 12\% in V2am), suggesting selective impairment of prediction error signaling rather than a general suppression of visual responses.

\subsection{Motor task activation by cell type}

To confirm that L5 IT neurons are not simply more responsive overall, we analyzed locomotion-onset responses across conditions (Table~\ref{tab:motor}).

\begin{table}[H]
\centering
\caption{V1 calcium response amplitude ($\Delta F/F$, \%) at locomotion onset in closed-loop and open-loop conditions by cell type. Paired $t$-test (closed vs.\ open within cell type). Hierarchical bootstrap 90\% CI.}
\label{tab:motor}
\begin{tabular}{llcccc}
\toprule
\textbf{Cre Line} & \textbf{Layer} & \textbf{Closed (\%)} & \textbf{Open (\%)} & \textbf{$\Delta$ (C$-$O)} & \textbf{90\% CI} \\
\midrule
Tlx3 & L5 IT & $3.42 \pm 0.38$ & $1.28 \pm 0.21$ & $\mathbf{2.14}$ & $[1.72, 2.56]$ \\
Ntsr1 & L6 & $2.87 \pm 0.41$ & $1.14 \pm 0.28$ & $1.73$ & $[1.18, 2.28]$ \\
SST & IN & $1.94 \pm 0.32$ & $0.98 \pm 0.19$ & $0.96$ & $[0.58, 1.34]$ \\
VIP & IN & $1.78 \pm 0.29$ & $1.02 \pm 0.22$ & $0.76$ & $[0.38, 1.14]$ \\
Scnn1a & L4 & $1.52 \pm 0.24$ & $1.21 \pm 0.20$ & $0.31$ & $[-0.02, 0.64]$ \\
Cux2 & L2/3 & $1.38 \pm 0.22$ & $1.14 \pm 0.18$ & $0.24$ & $[-0.06, 0.54]$ \\
\bottomrule
\end{tabular}
\end{table}

L5 IT neurons showed the largest closed-vs-open difference ($\Delta = 2.14\%$), confirming their role as the primary cortical substrate for visuomotor integration.

\section{Discussion}

Our findings reveal a striking layer specificity in the cortical effects of antipsychotic drugs. Three mechanistically distinct antipsychotics---spanning typical (haloperidol), atypical (clozapine), and partial agonist (aripiprazole) pharmacology---converge on reducing long-range correlations specifically in L5 IT neurons. This convergence suggests that decorrelation of deep-layer inter-areal communication represents a shared therapeutic mechanism independent of specific receptor profiles.

The preferential reduction of mismatch (prediction error) responses over sensory-driven grating responses aligns with computational theories of psychosis as a failure of predictive processing \citep{friston2005theory, adams2013computational}. In this framework, psychosis arises from aberrant precision-weighting of prediction errors. Antipsychotic-induced reduction of L5 IT prediction error responses and their inter-areal propagation could restore the balance between internally generated predictions and externally driven signals.

The absence of L5 IT decorrelation with amphetamine is notable. Amphetamine increases dopamine release and can precipitate psychotic symptoms, suggesting that its mechanism of action operates through different circuit elements than those affected by antipsychotic-induced decorrelation. This dissociation also argues against a simple dopaminergic explanation, as clozapine and aripiprazole have distinct D2 receptor pharmacology.

The layer specificity of antipsychotic effects is consistent with known cortical architecture. L5 IT neurons send long-range projections to other cortical areas and are the primary source of cortico-cortical communication \citep{harris2019hierarchical, gerfen2018long}. By selectively decorrelating L5 IT inter-areal interactions, antipsychotics may disrupt the propagation of aberrant internal models while preserving local sensory processing in superficial layers.

Limitations include the use of widefield imaging, which cannot resolve individual neurons, and the single-dose acute drug paradigm, which may not fully capture chronic treatment effects. The crystal skull preparation, while minimizing hemodynamic artifacts compared to clear skull approaches, still requires careful control \citep{kim2016crystal}. Future studies should examine chronic drug effects and use two-photon imaging for single-neuron resolution.

\section{Methods}

\subsection{Animals and transgenic lines}

All experiments were approved by the institutional animal care committee. Eight Cre lines were used: C57BL/6 with AAV-PHP.eB ($n = 6$), Emx1-Cre ($n = 4$), Cux2-CreERT2 ($n = 4$), Scnn1a-Cre ($n = 7$), Tlx3-Cre $\times$ Ai148 ($n = 25$), Ntsr1-Cre ($n = 3$), PV-Cre ($n = 2$), VIP-Cre ($n = 6$), and SST-Cre ($n = 5$) \citep{daigle2018suite, dana2014thy1}.

\subsection{Widefield calcium imaging}

Crystal skull cranial windows were implanted over dorsal cortex \citep{kim2016crystal}. GCaMP6 was excited at 470~nm LED illumination; images were acquired with a sCMOS camera at 30~Hz. Isosbestic imaging at 405~nm was used for hemodynamic correction.

\subsection{Behavioral paradigm}

Head-fixed mice ran on an air-supported spherical treadmill in a virtual corridor with vertical black and white bars. Closed-loop: locomotion drove visual flow. Open-loop: visual flow was replayed from a previous closed-loop session. Mismatch events were brief halts in visual flow during locomotion. Drifting gratings were presented during quiescence.

\subsection{Drug administration}

Single subcutaneous injections were administered during imaging sessions: clozapine (1~mg/kg), aripiprazole (3~mg/kg), haloperidol (0.1~mg/kg), or amphetamine (2~mg/kg). Pre-drug baseline was recorded for 30~min, followed by 60~min post-drug recording.

\subsection{Inter-regional correlation analysis}

Twelve bilateral dorsal cortical ROIs were defined. Pairwise Pearson correlations were computed between ROI time series in 5-min windows. Inter-regional distance was normalized by bregma-lambda distance. Pre-drug and post-drug correlation matrices were compared.

\subsection{Statistical analysis}

Hierarchical bootstrap with 90\% confidence intervals was used for all time-course and bar comparisons \citep{efron1993bootstrap}. Multiple comparison correction was applied where indicated. Effect sizes were computed as Cohen's $d$ for paired comparisons.

\section{Conclusion}

Antipsychotic drugs selectively decorrelate long-range cortical communication in layer 5 IT neurons while sparing superficial layers. This layer-specific mechanism, shared across three pharmacologically distinct antipsychotics but absent with a psychostimulant, provides a circuit-level account of antipsychotic action. The preferential reduction of prediction error responses and their inter-areal propagation connects these findings to predictive processing theories of psychosis, suggesting that restoring the balance between internal predictions and external signals may be a fundamental component of antipsychotic efficacy.

\bibliographystyle{plainnat}
\bibliography{references}

\end{document}
